---
format: yohaku/0.2
target: text
macro: '[[macros/raw]]'
---

## 正文
````text
"""Text addresses shared by Yohaku agents and the human JSON/CLI interface.

Positions are 1-based Unicode code-point line/column pairs. Ranges include
their start and exclude their end. LF and CRLF are preserved byte-for-byte;
neither half of CRLF is an addressable character. An empty final line exists
after a trailing newline. These are document coordinates, not screen pixels.
"""
import hashlib

MAX_RANGE_BYTES = 128 * 1024
MAX_TEXT_BYTES = 2 * 1024 * 1024


def lines(text):
    if not isinstance(text, str) or len(text.encode('utf-8')) > MAX_TEXT_BYTES:
        raise ValueError('Coordinate editing supports UTF-8 text up to 2 MiB')
    result = []
    offset = 0
    # Only LF/CRLF are line separators. Other Unicode separators stay in text.
    chunks = text.split('\n')
    for index, chunk in enumerate(chunks):
        newline = index < len(chunks)-1
        content = chunk[:-1] if newline and chunk.endswith('\r') else chunk
        result.append((offset, content))
        offset += len(chunk) + int(newline)
    return result


class RangeError(ValueError):
    """Out-of-range coordinate with program-computed document metrics."""
    def __init__(self, message, line_count, eof_line, eof_column):
        self.line_count = line_count
        self.eof_line = eof_line
        self.eof_column = eof_column
        super().__init__(
            f'{message}; document has {line_count} lines, last position is '
            f'{eof_line}:{eof_column}; read the range again or omit '
            f'end_line/end_column to page')


def metrics(table):
    return len(table), len(table[-1][1])+1


def position(table, line, column):
    if type(line) is not int or type(column) is not int or line < 1 or column < 1:
        raise ValueError('Line and column must be positive integers starting at 1')
    line_count, eof_column = metrics(table)
    if line > line_count or column > len(table[line-1][1])+1:
        raise RangeError(f'Text coordinate {line}:{column} is outside the document',
                         line_count, line_count, eof_column)
    return table[line-1][0]+column-1


def bounds(table, start_line, start_column, end_line, end_column):
    start = position(table, start_line, start_column)
    end = position(table, end_line, end_column)
    if end < start:
        line_count, eof_column = metrics(table)
        raise RangeError('Range end must not precede its start',
                         line_count, line_count, eof_column)
    return start, end


def read(text, start_line=1, start_column=1, end_line=None, end_column=None):
    table = lines(text)
    position(table, start_line, start_column)
    if end_line is None:
        end_line = min(start_line+79, len(table))
    line_count, eof_column = metrics(table)
    if type(end_line) is not int or not 1 <= end_line <= line_count:
        raise RangeError(f'End line {end_line} is outside the document',
                         line_count, line_count, eof_column)
    if end_column is None:
        # Default pages include their final newline when there is one. The
        # returned end can therefore be the next line's first column.
        if end_line < line_count:
            end_line, end_column = end_line+1, 1
        else:
            end_column = len(table[end_line-1][1])+1
    start, end = bounds(table, start_line, start_column, end_line, end_column)
    selected = text[start:end]
    if len(selected.encode('utf-8')) > MAX_RANGE_BYTES:
        raise ValueError('Requested range exceeds 128 KiB; request fewer lines or columns')
    has_more = end < len(text)
    return {
        'start_line':start_line, 'start_column':start_column,
        'end_line':end_line, 'end_column':end_column, 'text':selected,
        'line_count':line_count, 'eof_line':line_count,
        'eof_column':eof_column, 'has_more':has_more,
        'next':({'start_line':end_line,'start_column':end_column} if has_more else None),
        'coordinates':'1-based Unicode characters; start inclusive, end exclusive',
    }


def replace(text, start_line, start_column, end_line, end_column, expected, replacement):
    if not isinstance(expected,str) or not isinstance(replacement,str):
        raise ValueError('Expected text and replacement must be strings')
    table = lines(text)
    start, end = bounds(table, start_line, start_column, end_line, end_column)
    if text[start:end] != expected:
        raise ValueError('Selected text does not match expected; no edit made, read the range again')
    if len(expected.encode('utf-8')) > MAX_RANGE_BYTES or len(replacement.encode('utf-8')) > MAX_RANGE_BYTES:
        raise ValueError('A coordinate replacement must not exceed 128 KiB')
    result = text[:start]+replacement+text[end:]
    lines(result)  # Validate the resulting UTF-8 size before any write.
    return result


def verify(text, expected):
    if text != expected:
        raise ValueError('Read-back does not match the saved edit; inspect current text before retrying')
    return hashlib.sha256(text.encode('utf-8')).hexdigest()
````

## 输出与关联
```yaml
output: "agent/yohaku_coordinates.py"
```

## AI 笔记
```yaml
{"source_version": 1, "writer": "Yohaku release packaging", "summary": "Portable release source. Scope and packaging adjustments are documented in PROVENANCE.md."}
```
